\section{设计 Agent：从 brief 到交付}

导读把 Lovart 定位为垂直 Agent，本章进一步拆它的产品逻辑。设计不是单次生成图片，而是多轮工作流：理解 brief，收集素材，提出概念，生成方案，按反馈迭代，最后交付可用文件。设计师需要的不只是“生成一张图”，而是一个懂上下文、懂审美、能持续迭代的助手。

\begin{figure}[H]
\centering
\includegraphics[width=0.96\textwidth]{design-agent-workflow.png}
\caption{设计 Agent 工作流：从 brief 到素材、方案、迭代和交付。自制概念图，依据 00:52:20--00:57:00 对谈内容整理。}
\end{figure}

\begin{knowledgebox}{读图：设计工作流为什么适合垂直 Agent}
设计任务有明确上下文、强审美判断、多轮反馈和交付物要求。通用模型可以生成素材，但垂直 Agent 要负责把素材和用户目标串成连续工作流。
\end{knowledgebox}

\begin{importantbox}{设计 Agent 的本质}
设计 Agent 不是图片生成器，而是设计工作流协调器。它要把用户目标、品牌约束、素材、审美判断和交付格式放在同一个循环里。
\end{importantbox}

\subsection{设计师 AI Native 需求}

上一节讲工作流，本节具体看设计师为什么需要 AI Native 工具。设计师使用 AI，不只是为了省时间，还为了扩大探索空间。Lovart 需要处理品牌上下文、素材管理、版本迭代、协作和最终交付。它越像专业工作台，越不容易被一个通用聊天框完全替代。

\begin{figure}[H]
\centering
\includegraphics[width=0.96\textwidth]{designers-ai-native.png}
\caption{设计师 AI Native 需求：专业设计场景需要审美、上下文、版本和交付。自制概念图，依据 00:33:00--00:57:00 对谈内容整理。}
\end{figure}

\begin{knowledgebox}{术语消化：设计 Agent 的能力栈}
\begin{tabularx}{\textwidth}{p{0.22\textwidth}p{0.34\textwidth}X}
\toprule
能力 & 作用 & 难点 \\
\midrule
Brief 理解 & 把自然语言需求转成设计目标 & 用户表达常常模糊。 \\
素材上下文 & 管理品牌、图片、参考和限制 & 信息多且不断变化。 \\
审美判断 & 选择风格和视觉方向 & 主观性强，评价难。 \\
版本迭代 & 多轮修改和比较 & 需要记住历史意图。 \\
交付文件 & 输出可用设计资产 & 格式、协作和后续编辑很重要。 \\
\bottomrule
\end{tabularx}
\end{knowledgebox}

\begin{knowledgebox}{设计场景的深水区}
\begin{tabularx}{\textwidth}{p{0.22\textwidth}p{0.34\textwidth}X}
\toprule
深水区 & 具体表现 & 通用模型为什么不够 \\
\midrule
品牌一致性 & 色彩、字体、调性、资产规则 & 需要长期上下文和资产管理。 \\
多方案探索 & 同一 brief 下快速比较方向 & 需要版本管理和评审机制。 \\
团队协作 & 设计师、运营、客户、老板共同修改 & 需要评论、权限和历史记录。 \\
交付格式 & 可编辑文件、素材包、尺寸适配 & 需要进入真实生产工具链。 \\
\bottomrule
\end{tabularx}
\end{knowledgebox}

\begin{importantbox}{设计工作流的真正难点}
设计工作的难点不是生成一张“好看”的图，而是让一组视觉资产在品牌、目标用户、项目约束和团队反馈中保持一致。Lovart 如果要成为工作台，就必须管理这些连续上下文。
\end{importantbox}

\subsection{本章小结}

Lovart 的产品价值在于把 AI 生成能力放进设计工作流，而不是停留在一次性生成。垂直 Agent 的壁垒来自专业上下文和连续交付。

